\subsection{The symmetrizer coefficient from the same mixed action}

\begin{result}{E.180}{Proved}
{The full-GT mixed comparison identifies the historical coefficient}
Let \(F=(\lambda,f)\) be a full profinite GT element,
and let \(c_2(f)\) be the exact Schreier coefficient
of E.179.1.  Let \(g_F\) be the symmetrizer with
\[
 f(x,y)=g_F(y,x)^{-1}g_F(x,y),\qquad
 [g_F]_{\rm ab}=(\rho,\rho).
\tag{E.180.1}
\]
Then
\[
                         \rho=c_2(f)=-A_F(-1,1).
\tag{E.180.2}
\]
Here \(\rho\) is precisely the pure order-two frame
coefficient used in the historical CM character,
with its original sign.  Consequently the identity
of E.179 is relation \(({\rm IV})\) of
\citet[\S1, author p.3]{SymNakamuraSchneps2000}
with that paper's symmetrizer coefficient.
The proof uses the actual common-source comparison
E.176--E.179.  It does not infer the equality from
the agreement of the two coefficients on Galois
elements, and does not establish relation
\(({\rm III}')\) or the later signed CM covariance.
\end{result}

\begin{proof}
Throughout, \(\operatorname{Ad}(b)(z)=bzb^{-1}\).
All the group identifications below are made before
completion, and then on full profinite completions.
In the central quotient of
\(B_3=\langle u,w\mid uwu=wuw\rangle\), put
\[
 \begin{gathered}
 \mathfrak B=\widehat{C_2(s)*C_3(r)},\qquad
 s=uwu,\quad r=uw,\quad \tau=sr=u^{-1},\\
 x=\tau^{-2}=u^2,\qquad y=sxs=w^2,\qquad
 \mathfrak F=\ker(\mathfrak B\longrightarrow S_3)
                =\widehat F(x,y).
 \end{gathered}
\tag{E.180.3}
\]
For example \(u(sr)=(uw)^3\), and \(s^2=r^3=(uw)^3\);
these give the indicated relations in the cap.
The equalities for \(y\) follow from the same braid
relation.  Thus this is the literal modular marking
of the historical pure frame, not a new inner
normalization of it.

The standard full-GT braid action of E.179.3 induces
an automorphism \(\alpha\) of \(\mathfrak B\) with
\[
 \alpha(\tau)=\tau^\lambda,\qquad
 \alpha(x)=x^\lambda,\qquad
 \alpha(y)=f(x,y)^{-1}y^\lambda f(x,y).
\tag{E.180.4}
\]
It induces the identity on the label group \(S_3\).
There is a unique \(\beta\in\mathfrak F\) such that
\[
                         \alpha(s)=\beta s\beta^{-1}.
\tag{E.180.5}
\]
Indeed every order-two element in the full profinite
free product is conjugate to the \(C_2\) factor.
The label of a conjugator in E.180.5 centralizes the
named transposition, so multiplication on the right
by \(1\) or \(s\) makes it pure.  Its uniqueness
follows from
\(C_{\mathfrak B}(s)=\langle s\rangle\), whose
intersection with \(\mathfrak F\) is trivial.
These are precisely the free-factor statements of
\citet[Theorems A and B, p.155; proofs,
pp.157--160]{SymHerfortRibes1985}.
We use their full profinite free product theorem;
no corresponding centralizer claim for a maximal
pro-\(\ell\) quotient is being substituted.

We first identify the coefficient without relying
on historical notation.  Let
\(\vartheta=\operatorname{Ad}(s)|_{\mathfrak F}\);
then \(\vartheta(x)=y,\vartheta(y)=x\).
Computing \(\alpha(y)\) from \(y=sxs\) and
E.180.5, and comparing with E.180.4, gives
\[
              f\,\beta\,\vartheta(\beta)^{-1}
                          \in C_{\mathfrak F}(y^\lambda)
                           =\langle y\rangle^\wedge .
\tag{E.180.6}
\]
The last equality is again the free-factor centralizer
theorem, now in
\(\mathfrak F=\langle x\rangle^\wedge
 \amalg\langle y\rangle^\wedge\).
Write \([\beta]_{\rm ab}=(b_1,b_2)\).
Since \(f\) has zero abelianization, E.180.6 gives
\[
                  (b_1-b_2,b_2-b_1)=(0,t)
\]
for some \(t\in\widehat{\mathbb Z}\).  Therefore
\(b_1=b_2\), \(t=0\), and
\[
 f=\vartheta(\beta)\beta^{-1},\qquad
 g_F=\beta^{-1},\qquad
 [\beta]_{\rm ab}=(-\rho,-\rho).
\tag{E.180.7}
\]
The uniqueness of this symmetrizer follows in the
same way from the centralizer of \(s\) in
\(\mathfrak F\rtimes\langle s\rangle\).
This agrees exactly with
\citet[Definition 5.1, author p.17]
{SymNakamuraSchneps2000} and the historical convention
\((\beta^{-1})^{\rm ab}=(\rho,\rho)\).

It remains to connect this \(\beta\) to the specific
mixed action used in E.179.  This step retains the
change of marking in the second input of \(f\).
Let
\[
 \begin{gathered}
 R_H=\langle x,s\rangle\leq\mathfrak B,\qquad
 R=\widehat{\langle v,d\mid(vd)^2=1\rangle},\\
 j:R\xrightarrow{\;\sim\;}R_H,\qquad
 j(v)=x,\quad j(d)=x^{-1}s,\quad
 j(d^2)=x^{-1}y^{-1}.
 \end{gathered}
\tag{E.180.8}
\]
This is the completion of the actual index-three
subgroup: it is the inverse image of the subgroup
generated by the named transposition in \(S_3\).
Its index-two pure subgroup has basis \(x,y\).
In particular its induced topology is the full
profinite topology.

Put \(m=(\lambda-1)/2\) and
\(\theta=\operatorname{Ad}(\tau)|_{\mathfrak F}\).
The literal words give
\[
 \theta(x)=x,\qquad
 \theta(y)=x^{-1}y^{-1}=j(d^2),\qquad
 \alpha\theta=\operatorname{Ad}(x^{-m})\theta\alpha.
\tag{E.180.9}
\]
The last equality uses
\(\alpha(\tau)\tau^{-1}=\tau^{\lambda-1}=x^{-m}\).
Consequently \(\operatorname{Ad}(x^m)\alpha|_{R_H}\)
acts on \(x\) by \(x^\lambda\), and on \(j(d^2)\)
by
\[
 \theta(f)^{-1}j(d^2)^\lambda\theta(f),\qquad
                 \theta(f)=f(x,x^{-1}y^{-1}).
\tag{E.180.10}
\]
Under \(j\), this last word is
\(h=f(v,d^2)\), exactly as in E.179.4.
Thus the action in E.180.10 and \(j\phi j^{-1}\)
agree on \(v\) and \(d^2\).
Their images of \(d\) agree as well: after conjugation
by \(h\), both have square \(d^{2\lambda}\), and
therefore centralize that nontrivial element of
the free factor \(\langle d\rangle^\wedge\) in
\[
                  R=\langle a\mid a^2=1\rangle
                            \amalg\langle d\rangle^\wedge,
                   \qquad a=vd .
\]
The free-factor centralizer theorem places both
images in \(\langle d\rangle^\wedge\).
Multiplication by two is injective on
\(\widehat{\mathbb Z}\), so their squares determine
them uniquely.  We have proved the exact formula
\[
             j\phi j^{-1}
                  =\operatorname{Ad}(x^m)\alpha|_{R_H}.
\tag{E.180.11}
\]
In particular, if we identify \(R_H\) with \(R\)
and put \(b=v^m\beta\), then
\[
                             \phi(a)=bab^{-1}.
\tag{E.180.12}
\]
No replacement of \(f(x,x^{-1}y^{-1})\) by \(f(x,y)\)
has occurred.

Write \(c=c_2(f)\), \(k=f(d,v^2)\), and
\[
                     g=d^{-4c}h,\qquad S(v)=d,\quad S(d)=v .
\tag{E.180.13}
\]
E.179 gives \(kg=v^{4c}\).  Checking on \(v,d\)
therefore gives the literal equality
\[
                          S\phi S=\operatorname{Ad}(g)\phi .
\tag{E.180.14}
\]
Use \(S(a)=da d^{-1}\), E.180.12, and
\(S(\phi(d))=k^{-1}v^\lambda k\).
Computing the image of \(a\) under the two sides
of E.180.14 yields two conjugators
\[
                        gb,\qquad
                 z=k^{-1}v^\lambda k\,S(b)d .
\tag{E.180.15}
\]
Their quotient centralizes \(a\).  By the same
free-factor theorem \(C_R(a)=\langle a\rangle=C_2\).
It suffices to apply the single continuous character
\[
  L:R\longrightarrow\widehat{\mathbb Z},\qquad
                  L(v)=1,\quad L(d)=-1,\quad L(a)=0 .
\tag{E.180.16}
\]
We have \(L\circ S=-L\).  Both \(x=v\) and
\(y=ava\) have \(L\)-value one, so E.180.7 gives
\[
 L(\beta)=-2\rho,\quad L(b)=m-2\rho,\quad
 L(h)=L(k)=0,\quad L(g)=4c .
\tag{E.180.17}
\]
The quotient of the two conjugators in E.180.15
has \(L\)-value zero.  Hence
\[
 \begin{aligned}
  4c+m-2\rho
       &=L(gb)=L(z)\\
       &=\lambda-m+2\rho-1,\\
  4(c-\rho)&=\lambda-1-2m=0 .
 \end{aligned}
\tag{E.180.18}
\]
The group \(\widehat{\mathbb Z}=\prod_p\mathbb Z_p\)
is torsion free, so \(c=\rho\).
The remaining equality in E.180.2 is the exact
Schreier--Alexander calculation in E.179.

Nakamura--Schneps Proposition 5.3 proves the same
coefficient equality under their relation
\(({\rm III}')\).  Here it is obtained instead
from the independently proved common-source
mixed comparison and the explicit modular
marking E.180.8--E.180.11.
Only this equality and relation \(({\rm IV})\)
are concluded.  A theorem that also assumes
\(({\rm III}')\), a labelled cover comparison,
or a signed genus-one covariance square still
requires that additional input.
\end{proof}
